\section{The actual rational full matrix: exterior and boundary necessity}
\label{sec:polar-exterior-obstruction}

The rational completion in MA-MO-016 carries actual polar source
sections. A compactly supported functional of either represented
component is therefore bounded on that completion, although a bare
polar exponential does not lie in its diagonal Hilbert summand.
This observation transfers the actual-source obstruction to the
principal character without discarding either polar coordinate.

Keep \(S=\{\infty\}\), the arbitrary smooth real cutoff \(\phi\)
of MA-MO-018, and the original common source maps. Put
\[
 \mathcal W_{a,s}=H_{\Sigma,a,s}\oplus\mathbb C^2,
 \qquad 0<a<1/2,\quad s\geq0,\qquad
 \mathscr D_g^{\rm raw}=C_g(P_+-P_-).
 \tag{MA-MO.650}
\]
The raw map here is defined on \(\mathscr S_{\rm sum}\), with
its actual pair realization \(\mathfrak s\) from (MA-MO.252).
No zero-free half-plane is assumed in this section.

\paragraph{MA-MO-026 (necessary full-divisor retention in the polar completion).}
Let \(\rho\) be a nontrivial zeta zero of multiplicity \(m\).
For every nonzero \(g\in C_c^\infty(\mathbb R)\), the original
map \(\mathscr D_g^{\rm raw}\) has no bounded endomorphism
extension on \(\mathcal W_{a,s}\) if either
\[
 \begin{array}{ll}
 \text{(a)}&|\Re\rho-1/2|>a,\\
 \text{(b)}&|\Re\rho-1/2|=a\quad\text{and}\quad s\leq2m-1.
 \end{array}
 \tag{MA-MO.651}
\]
This includes tests satisfying
\(g(0)=\mathcal Mg(1/2)=\mathcal Mg(-1/2)=0\), and tests
whose Mellin transform vanishes at \(\rho-1/2\). In case (a)
there is a fixed compact smooth negative source \(i_-h\) and
\(t_n\to+\infty\) such that
\[
 \frac{\|\mathscr D_g^{\rm raw}i_-U_{t_n}h\|_{\mathcal W_{a,s}}}
      {\|i_-U_{t_n}h\|_{\mathcal W_{a,s}}}\longrightarrow\infty.
 \tag{MA-MO.652}
\]
Every displayed input and output belongs to the original smooth
common source. In case (b) the contradiction uses convergent
integrals of these same actual inputs, and compact smooth
approximations to those integrals.

\begin{proof}
\emph{The bounded output functional and the literal rational source.}
For \(v\in C_c^\infty(\mathbb R)\) define
\[
 \Lambda_v(u,c)=\int v(x)
       \{u(x)+c_0\Theta f_0(x)+c_1\Theta f_1(x)\}\,dx.
 \tag{MA-MO.653}
\]
This is the bilinear integral of \(v\) against the first actual
pair component of \(\mathfrak s(u,c)\). It is bounded on
\(\mathcal W_{a,s}\): the diagonal term is bounded by weighted
Cauchy--Schwarz, and the two fixed smooth section functions are
integrable on the compact support of \(v\). The polar terms
are included in this functional, not subtracted from the source.
Also \(i_-y=(y,0)\) is an isometry from \(H_{\Sigma,a,s}\)
into \(\mathcal W_{a,s}\). These statements follow directly
from the actual completion, without its conditional matrix theorem.

For a compact smooth \(h\), let
\[
 F(x)=\sum_{n\geq1}\mu(n)n^{-1/2}h(x+\log n),
 \qquad A(w)=\zeta(1/2+w),\qquad q=1-\phi.
 \tag{MA-MO.654}
\]
This is the original wide inverse. It is locally a finite sum,
has a finite right support endpoint, and obeys the absolute
wide bounds (MA-MO.560), with \(\zeta(b+1/2)\) in place of
the ideal majorant. The actual additive source for \(i_-y\)
is (MA-MO.325--326), namely
\[
 u_y(t)=\frac{q(\log|t|)F_y(\log|t|)}{2\sqrt{|t|}},
 \qquad h_y=u_y+\mathfrak F u_{Ry}.
 \tag{MA-MO.655}
\]
For compact \(y\), \(qF_y\) has compact logarithmic support,
so \(u_y\) is an actual compact annular even Schwartz function.
The source formula (MA-MO.655) consequently holds before any
narrow contour shift. The conditional estimates in MA-MO-018
are not being imported here. For general \(y\in\mathcal E\),
the arbitrarily rapid right tail and the literal wide estimates
still make \(u_y\) Schwartz and zero near additive zero;
compact right support is not needed. Thus the original maps
act on their stated smooth common source without a narrow
boundedness premise.

The normalization of the actual theta map gives exactly
\[
 \Theta u_y(x)=\sum_{n\geq1}n^{-1/2}
                   q(x+\log n)F_y(x+\log n).
 \tag{MA-MO.656}
\]
If \(y=U_th\) and \(t\) is sufficiently large, the support
argument (MA-MO.561--562) gives \(qF_{RU_th}=0\).
Hence \(h_{U_th}=u_{U_th}\) exactly, before Fourier
transformation. Since \(P_-i_-y=i_-y\), the first component
of the actual raw output is \(C_g(\Theta u_{U_th}-U_th)\).
With \(v_g=C_{g^\vee}v\) and
\[
 W(y)=\sum_{n\geq1}n^{-1/2}v_g(y-\log n),
 \qquad f(t)=\Lambda_v(\mathscr D_g^{\rm raw}i_-U_th),
\]
the literal finite-support calculation therefore gives
\[
                 f(t)=\int q(y)W(y)F(y-t)\,dy
                         \qquad(t\geq t_0).
 \tag{MA-MO.657}
\]
The subtracted input pairs to zero after increasing \(t_0\).
This establishes the exact hypotheses of the paired calculation
in MA-MO-024 on the actual polar completion.

\emph{The pole-cleared Laplace identity.}
Use the definitions \(H_h,V_g,B_v,\mathcal R_{h,v}\) of
(MA-MO.558,.564), now with the ordinary positive integers.
The proof (MA-MO.559--567) applies verbatim to these literal
sums: both \(B_v\) and \(\mathcal R_{h,v}\) are entire,
the latter is a compact double integral, and for \(\Re w>1/2\)
the absolutely convergent actual paired Laplace transform obeys
\[
 \mathfrak F(w)=H_h(w)V_g(-w)
             -\frac{H_h(w)B_v(w)}{A(w)}-\mathcal R_{h,v}(w).
 \tag{MA-MO.658}
\]
The only change from the nontrivial-character case is the pole
of \(A\) at \(w=1/2\). Clear it by the entire function
\(\widetilde A(w)=(w-1/2)A(w)\). Thus a bounded extension
of the raw matrix, which would force \(\mathfrak F\) to be
holomorphic on \(\Re w>a\), would imply
\[
 \widetilde A(w)
  [H_h(w)V_g(-w)-\mathcal R_{h,v}(w)-\mathfrak F(w)]
       =(w-1/2)H_h(w)B_v(w),\qquad \Re w>a.
 \tag{MA-MO.659}
\]
This follows by the identity theorem from the original wide
half-plane. Clearing the zeta pole does not change the order
of any nontrivial zero \(\lambda=\rho-1/2\), since
\(\lambda\ne1/2\).

At any such \(\lambda\) the nonannihilation construction
(MA-MO.568--572) also applies without change: its coefficients
are now \(n^{-1/2-\lambda}\), the term \(n=1\) is the sole
term near the last support edge, and the majorant is the
convergent ordinary zeta series. For every complex nonzero
\(g\) it therefore supplies a compact smooth \(v\) with
\(B_v(\lambda)\ne0\). Choose
\(h(x)=e^{-\lambda x}h_0(x)\) for a nonnegative nonzero
compact smooth \(h_0\), so \(H_h(\lambda)>0\).
No nonzero value of \(\mathcal Mg(\lambda)\) is required.

If \(\Re\lambda>a\), (MA-MO.659) is already a contradiction
at \(w=\lambda\). More precisely the argument of
(MA-MO.573--575), using the bounded functional (MA-MO.653)
and the isometric negative embedding, proves (MA-MO.652).
Every raw output is still a smooth common-source vector by
(MA-MO.655). This proves case (a) for a right exterior zero.

\emph{A boundary zero and its exact polynomial threshold.}
Suppose \(\lambda=a+i\gamma\) has multiplicity \(m\), and
assume a bounded extension exists. For \(\delta>0\), the
Bochner integral
\[
 y_\delta=\int_{t_0}^{\infty}
                   e^{-(\lambda+\delta)t}U_th\,dt
                    \quad\text{in }H_{\Sigma,a,s}
 \tag{MA-MO.660}
\]
converges because the translation norm is at most a constant
times \(e^{at}(1+t)^{s/2}\). Truncating the integral at a
finite upper endpoint gives a compact smooth original source;
the truncations converge in \(H\), and their negative embeddings
converge in \(\mathcal W\). Boundedness then justifies
\[
 \Lambda_v(\mathscr D_g i_-y_\delta)
                        =\mathfrak F(\lambda+\delta).
 \tag{MA-MO.661}
\]
If \(\operatorname{supp}h\subset[\alpha,\beta]\), direct
substitution in (MA-MO.660) gives
\[
 y_\delta(x)=e^{-(\lambda+\delta)x}
       \int_{-\infty}^{x-t_0}e^{(\lambda+\delta)u}h(u)\,du.
 \tag{MA-MO.662}
\]
It vanishes for \(x<\alpha+t_0\), is uniformly bounded with
all derivatives on each compact set as \(0<\delta\leq1\),
and for \(x\geq\beta+t_0\) equals
\(H_h(\lambda+\delta)e^{-(\lambda+\delta)x}\).
Since the exact sum weight is
\((e^{2ax}+e^{-2ax})(1+x^2)^{s/2}\), integration of this
tail gives
\[
                  \|y_\delta\|_H
                         =O(\delta^{-(s+1)/2}).
 \tag{MA-MO.663}
\]
Increase \(t_0\), if necessary, so that \(\beta+t_0\geq1\).
The contribution on the remaining fixed compact interval is
uniformly bounded. The tail bound follows, for example, by
\(u=2\delta x\) in the integral of
\(e^{-2\delta x}(1+x)^s\), with a harmless fixed lower
endpoint. The other exponential in the sum weight only improves
integrability.

If \(s<2m-1\), it follows that \(\delta^m y_\delta\to0\)
strongly in \(H\). If \(s=2m-1\), this family is bounded
and tends to zero in \(L^2\) on every compact set. Pairing
first with compactly supported \(H\)-vectors and then using
their density and the uniform bound shows
\(\delta^m y_\delta\rightharpoonup0\) weakly in \(H\).
In either case a fixed bounded functional composed with the
assumed bounded matrix gives
\[
       \delta^m\Lambda_v(\mathscr D_g i_-y_\delta)\longrightarrow0.
 \tag{MA-MO.664}
\]
But near \(\lambda\) the right side of (MA-MO.658), justified
there by (MA-MO.659), has a pole of order exactly \(m\): its
numerator \(H_h(\lambda)B_v(\lambda)\) is nonzero, and
\(\mathcal R_{h,v},H_h,V_g\) are entire. If
\(A(w)=(w-\lambda)^m a_0(w)\) with \(a_0(\lambda)\ne0\),
then (MA-MO.661) instead gives
\[
 \delta^m\Lambda_v(\mathscr D_g i_-y_\delta)
       \longrightarrow
       -\frac{H_h(\lambda)B_v(\lambda)}{a_0(\lambda)}\ne0.
 \tag{MA-MO.665}
\]
This contradicts (MA-MO.664), proving the boundary assertion.

Finally the zeta functional equation carries a left exterior
or left boundary zero to a right one with the same multiplicity.
Applying the proved right-side argument to that actual zero
proves all of (MA-MO.651). This use of the functional equation
asserts no new location of zeros. The construction of \(v\)
only required \(g\ne0\), which proves the stated test scopes.
\end{proof}

\paragraph{The exact arithmetic meaning of all widths.}
For this actual completion, the following are equivalent:
zeta RH; existence, at each \(0<a<1/2\), of some \(s\geq0\)
and some nonzero compact smooth test for which the full original
matrix extends boundedly; and trace-class extension for every
such width, weight and test. Necessity is (MA-MO.651)(a) at
arbitrarily small widths. Sufficiency is precisely MA-MO-018
with \(a_*=0\); it explicitly uses zeta RH. The test may also
be required to satisfy \(g(0)=\mathcal Mg(\pm1/2)=0\), since
the nonzero class constructed after (MA-MO.577) has those
properties. Under the user-specified theorem A, this arithmetic
premise is supplied. It remains a named zero-location input in
an audit of the method.

At a fixed width, bounded full-source action therefore requires
no exterior zero and retention of every boundary multiplicity
under the sharp MA-WP condition \(s>2m-1\). Sufficiency without
a strictly smaller containing strip is not proved here. The
result concerns the actual polar completion, not the full
coordinate product already excluded by MA-MO-005, nor the
restricted positive source column of MA-MO-016/023.

The source locators are Meyer (2005),
\texttt{lem:Hilm\_p\_embed} and its original \(p_\pm\) formulas,
supplied \texttt{main.tex} lines 3808--4043, and the actual
completion and source calculations (MA-MO.250--263,.323--326).
In Grothendieck's cote~1 framework, scans 079--080 and 111--112
through BR-TR-01/02, a trace or Fredholm determinant requires a
defined bounded nuclear realization. Adjoining the two actual
polar sections does not remove the analytic obstruction from an
unretained zero or jet. The finite-length \(K_0\) class records
that multiplicity; it does not supply the missing operator norm.
No faithful full-jet positivity, motivic purity, or independent
arithmetic vanishing theorem is asserted.
